\section{Limitations}

While this work provides quantitative scaling principles for agent systems across architectures and model families, several limitations remain. 

\textbf{(i)} Our framework systematically compares canonical coordination structures with preliminary exploration of scaling number of agents up to nine. However, our empirical findings suggest that scaling to larger collectives may face fundamental barriers: the communication overhead we measured grows superlinearly with agent count, and coordination efficiency degrades substantially beyond moderate team sizes. Whether such collectives can exhibit beneficial emergent behaviors, such as spontaneous specialization or hierarchical self-organization, or whether communication bottlenecks dominate remains an open question that parallels phase transitions in complex adaptive systems. 

\textbf{(ii)} While we explore capability heterogeneity by mixing models of different intelligence levels within the same LLM family, all agents share identical base architectures differing only in scale and role prompts. \textcolor{black}{A preliminary investigation of 13 heterogeneous configurations on BrowseComp-Plus (mixing models within and across families) finds no evidence that model mixing bypasses the capability-saturation threshold: centralized heterogeneous configurations underperform their strong-model homogeneous counterparts by a mean of 12.6 percentage points, while decentralized configurations show marginal gains (+2.0 pp) largely attributable to the stronger constituent model (Table~\ref{tab:s2_heterogeneous}).} Future work should investigate teams combining different model architectures, domain-specialized fine-tuning, or complementary reasoning strategies to understand when \emph{epistemic diversity} yields robustness rather than coordination noise. Additionally, our heterogeneity experiments (Figure~\ref{fig:heterogeneous}) hint that certain models may be better suited for orchestration versus execution roles; systematic study of \emph{role-specialized} training or selection could enable more principled team composition. 

\textbf{(iii)} Our analysis reveals that tool-heavy environments represent a primary failure mode for multi-agent coordination, with significant negative interactions between tool count and system efficiency. Developing specialized coordination protocols for tool-intensive tasks, such as explicit tool-access scheduling, capability-aware task routing, or hierarchical tool delegation, represents an important direction for improving multi-agent reliability. 

\textbf{(iv)} While we controlled prompts to be identical across conditions for experimental validity, we did not optimize prompts specifically for each model or model family. Given known sensitivity of LLM behavior to prompt formulation, architecture-specific prompt tuning may yield different scaling characteristics than those reported here. 

\textcolor{black}{\textbf{(v)} Our analysis spans six agentic benchmarks, which, while diverse in task structure (deterministic tool use, quantitative reasoning, sequential planning, dynamic web navigation, software engineering, and CLI tasks), may not capture the full spectrum of agentic task characteristics. The strong differentiation in MAS effectiveness across these benchmarks (Figure \ref{fig:boxplots}) suggests that additional environments, particularly those with novel task structures such as embodied agents, multi-user interaction, or long-horizon temporal dependencies, would further strengthen confidence in the identified thresholds and scaling principles. SWE-bench Verified and Terminal-Bench use 20-instance subsets (smaller than the 50--100 instances used for BrowseComp-Plus, Finance Agent, PlanCraft, and Workbench) due to the computational cost of Docker-based evaluation environments. Bootstrap 95\% confidence intervals (10,000 resamples) yield typical widths of $\pm$20 percentage points per cell; while individual pairwise comparisons are underpowered at $n=20$, aggregate trends across 8 models $\times$ 2 benchmarks remain robust (e.g., the 45\% threshold achieves 94\% match rate, $p < 0.001$ by binomial test). All per-cell bootstrap CIs are reported in Table~\ref{tab:s6_bootstrap}.} 

\textbf{(vi)} The economic viability of multi-agent scaling remains a practical barrier, rooted in part in the token-centric communication paradigm: current coordination requires agents to serialize reasoning into natural language tokens (or at minimum, read shared context and output agreement signals), imposing fundamental latency and cost floors. Emerging approaches such as latent-space reasoning or direct activation sharing between models could circumvent this bottleneck, potentially altering the scaling dynamics we observe if inter-agent communication shifts from token exchange to more efficient representational transfer. As shown in our cost analysis (Section \ref{subsec:coord_eff}), token consumption and latency grow substantially with agent count, often without proportional performance gains. Future work should explore efficiency-oriented designs, such as sparse communication, early-exit mechanisms, or distilled coordinator models, to make multi-agent deployments economically feasible at scale. Complementary \emph{latency-oriented} designs, where parallel agent branches execute speculatively and suboptimal trajectories are pruned post-hoc, may trade increased total compute for reduced wall-clock time, a trade-off increasingly relevant for real-time applications where response latency dominates cost considerations. Additionally, current agentic benchmarks capture dynamic text-based environments but do not yet include long-horizon temporal dependencies or real-world feedback loops. Integrating embodied or multimodal settings (e.g., robotic control, medical triage, multi-user social interaction) will test whether our observed scaling principles generalize beyond symbolic domains.

\textcolor{black}{\textbf{(vii)} Our regression analysis clusters observations at the dataset level ($G=6$). With a small number of clusters, cluster-robust standard errors are known to be conservative, and several predictors that are significant under naive OLS lose significance under cluster-robust inference (Table~\ref{tab:s4_cluster}). We therefore report both naive and cluster-robust estimates throughout, framing dataset-level predictors as descriptive patterns supported by directional consistency rather than confirmed at conventional significance levels. Leave-one-dataset-out cross-validation further highlights the challenge of predicting \emph{absolute} success rates across structurally diverse domains without dataset-specific parameters; however, the within-domain cross-validated model correctly selects the optimal architecture in 87\% of held-out cases, indicating that relative coordination patterns transfer across domains even when absolute performance levels vary.}